\documentclass[compress]{beamer}
\usecolortheme{default}

\usepackage{lmodern, epstopdf, multirow, array, setspace, bbm, tcolorbox, pbox, bm, mathtools, tikz}
\usepackage[T1]{fontenc}

\setbeamertemplate{frametitle}[default][center]
\setbeamertemplate{footline}{\hfill \insertframenumber/\inserttotalframenumber\hspace*{5pt}\vspace{5pt}}
\setbeamertemplate{navigation symbols}{}

\newlength{\wideitemsep}
\setlength{\wideitemsep}{\itemsep}
\addtolength{\wideitemsep}{4pt}
\let\olditem\item
\renewcommand{\item}{\setlength{\itemsep}{\wideitemsep}\olditem}
\setlength{\leftmargini}{12pt}
\setlength{\leftmarginii}{14pt}
\setbeamertemplate{itemize/enumerate subbody begin}{\vspace{6pt}}
\setbeamertemplate{itemize item}[circle]
\setbeamertemplate{enumerate item}[circle]
\setbeamertemplate{enumerate subitem}[default]

\setbeamercolor{alerted text}{fg=blue}
\setbeamersize{text margin left=1.5em,text margin right=1.5em}

\tcbset{
    colback=white,
    colframe=black,
    coltext=black,
    coltitle=blue,
    colbacktitle=white,
    boxrule=1pt,
    titlerule=0pt,
    toptitle=4pt,
    top=2pt
}

\newenvironment{changemargin}[2]{% 
  \begin{list}{}{% 
    \setlength{\topsep}{0pt}% 
    \setlength{\leftmargin}{#1}% 
    \setlength{\rightmargin}{#2}% 
    \setlength{\listparindent}{\parindent}% 
    \setlength{\itemindent}{\parindent}% 
    \setlength{\parsep}{\parskip}% 
  }% 
  \item[]}{\end{list}}
  
\newenvironment{defn}{%
  \begin{tcolorbox}[title={Definition}]%
  }{\end{tcolorbox}}
  
\newenvironment{prop}[1]{%
  \begin{tcolorbox}[title={Proposition #1}]%
  }{\end{tcolorbox}}
  
\newcommand{\cit}[1]{{\footnotesize \textcolor[rgb]{0.5,0,0}{#1}}}

% Font size for code listings (verbatim blocks); reset with \normalsize
\newcommand{\codesize}{\fontsize{7}{8.5}\selectfont}

\newcommand{\slidesetup}[1]{%
  \title{#1}%
  \subtitle{\texorpdfstring{\vspace{0.3\baselineskip}}{} {\small This version: \today}}%
  \author[Mikkel Plagborg-Moller]{ECMA 31150, Spring 2026 \texorpdfstring{\\[0.5ex]}{} University of Chicago \texorpdfstring{\\[0.5ex]}{} Department of Economics \texorpdfstring{\\[4ex]}{} Mikkel Plagborg-Moller}%
  \date{}%
}

% \circleoverlay[options]{text}
% Draws a red circle with a text box nearby on a tikz overlay.
% Options (with defaults):
%   cx=0cm, cy=0cm      - circle center (shift from page center)
%   r=1cm               - circle radius
%   tx=0cm, ty=0cm      - text box offset from circle center
%   text anchor=east    - anchor of text node (east = text left, west = text right)
%   font=\normalsize    - font size for text
\pgfkeys{
  /circleoverlay/.cd,
  cx/.store in=\cocx, cx=0cm,
  cy/.store in=\cocy, cy=0cm,
  r/.store in=\cor, r=1cm,
  tx/.store in=\cotx, tx=0cm,
  ty/.store in=\coty, ty=0cm,
  text anchor/.store in=\cotextanchor, text anchor=east,
  font/.store in=\cofont, font=\normalsize,
}
\newcommand{\circleoverlay}[2][]{%
  \pgfkeys{/circleoverlay/.cd, cx=0cm, cy=0cm, r=1cm, tx=0cm, ty=0cm, text anchor=east, font=\normalsize, #1}%
  \pgfmathsetlengthmacro{\circlediameter}{2*\cor}%
  \node[draw=red, line width=2pt, circle, minimum size=\circlediameter, inner sep=0pt]
    (tmpcirc) at ([xshift=\cocx,yshift=\cocy]current page.center) {};%
  \node[anchor=\cotextanchor, fill=white, draw=red, text=red, inner sep=4pt, align=left, font=\cofont]
    at ([xshift=\cotx,yshift=\coty]tmpcirc.center) {#2};%
}

\DeclareMathOperator*{\argmax}{argmax}
\DeclareMathOperator*{\argmin}{argmin}
\DeclareMathOperator{\diag}{diag}
\DeclareMathOperator{\var}{Var}
\DeclareMathOperator{\cov}{Cov}
\DeclareMathOperator{\corr}{Corr}
\DeclareMathOperator{\lrv}{LRV}
